\begin{table*}[t]
\centering
\label{tab:dataset_examples}
\setlength{\tabcolsep}{6pt}
\resizebox{\textwidth}{!}{%
\begin{tabular}{cccc}
\toprule
{\textbf{Image}} & {\textbf{PixMo QA}} & \textbf{Understanding QA} &{\textbf{Generation Prompt}} \\
\midrule
\addlinespace[4pt]
\raisebox{-0.5\height}{\includegraphics[width=5cm]{example-image-a}} &
\parbox[c]{4cm}{\raggedright\small A blue-toned photograph of a coastal scene with waves.} &
\parbox[c]{4cm}{\raggedright\small Describe the scene.} &
\parbox[c]{6cm}{\raggedright\scriptsize You are given a photograph of a coastal landscape. Please provide a comprehensive description of all visible elements including the water conditions, rock formations, sky color, and any atmospheric effects.} \\
\addlinespace[10pt]
\raisebox{-0.5\height}{\includegraphics[width=5cm]{example-image-a}} &
\parbox[c]{4cm}{\raggedright\small A warm sunset over a desert canyon with sandstone.} &
\parbox[c]{4cm}{\raggedright\small What colors are present?} &
\parbox[c]{6cm}{\raggedright\scriptsize Analyze the following desert landscape image. Identify all color gradients, geological textures, and lighting conditions visible in the scene.} \\
\addlinespace[10pt]
\raisebox{-0.5\height}{\includegraphics[width=5cm]{example-image-a}} &
\parbox[c]{4cm }{\raggedright\small A lush green forest with sunlight filtering through.} &
\parbox[c]{4cm }{\raggedright\small Identify vegetation types.} &
\parbox[c]{6cm}{\raggedright\scriptsize Examine the provided forest photograph and enumerate all distinguishable plant species or vegetation categories. Note the distribution of light and shadow.} \\
\addlinespace[10pt]
\raisebox{-0.5\height}{\includegraphics[width=5cm]{example-image-a}} &
\parbox[c]{4cm }{\raggedright\small A purple-tinted cityscape at twilight with neon.} &
\parbox[c]{4cm }{\raggedright\small Describe the environment.} &
\parbox[c]{6cm}{\raggedright\scriptsize You are presented with a nighttime urban scene. Provide a thorough description covering all visible architectural elements, signage, lighting sources, and human activity.} \\
\addlinespace[4pt]
\bottomrule
\end{tabular}
}% end resizebox
\caption{Examples of counting dataset.}
\end{table*}